\slajdid{09_1-s018}
\begin{frame}[label=09_1-s018]
Za određivanje $V_{IL}$ i $V_{IH}$ situacija nije tako vidljiva kao što će biti kod logičkih kola sa bipolarnim tranzistorima. Ovde ćemo to morati da uradimo preko definicije. Da nađemo tačke u kojima je pojačanje po apsolutnoj vrednosti jednako 1.\par\medskip
$V_{IL}$ će se naći očigledno u drugoj oblasti za koju važi
\[\frac{V_{DD}-V_O}{R_L}=\frac{k_n}{2}\frac1{1+\frac{(V_I-V_{Tn})}{L_nE_{Cn}}}(V_I-V_{Tn})^2\]
Gledajući kako se ponaša tranzistor za $V_I=V_{Tn}+\varepsilon$ možemo zaključiti da će on i za malo $\varepsilon$ ući u režim velikih pojačanja odnosno da će $(V_{IL}-V_{Tn})\ll L_nE_{Cn}$, pa taj član možemo zanemariti.\par\medskip
\centering Za određivanje $V_{IL}$
\[\frac{V_{DD}-V_O}{R_L}\approx\frac{k_n}{2}(V_I-V_{Tn})^2\]
\end{frame}
